\chapter{Why the Past Matters Less}
\index{stability!philosophical interpretation}

\section{A Recurring Question}

Throughout this book a single question has reappeared in every domain: does historical influence decay with interaction order? In learning systems it appeared as derivative decay. In semantic systems as branch attenuation. In physical fields as accessibility contraction. In memory as the forgetting curve. In institutions as precedential half-lives. In cosmology as accessible volume decay.

This recurrence is not an accident of the framework's design. Stability is a structural property of historical systems themselves.

\section{Three Faces of Stability}

\textbf{Stability as forgetting.} In memory systems, stability appears as forgetting---the progressive reduction of sensitivity to remote experiential interactions. A world without forgetting would be a world where nothing could be compressed, summarized, or reconstructed.

\textbf{Stability as abstraction.} In semantic systems, higher-order branch interactions are absorbed into coarser conceptual structures. Abstraction is the cognitive manifestation of stability.

\textbf{Stability as relaxation.} In physical systems, accessible configuration volume contracts as field dynamics evolve. The system loses sensitivity to increasingly remote alternative initial conditions.

\section{Stability as Historical Curvature}

A stable trajectory space is one in which nearby alternative histories \emph{converge under projection}: as truncation depth $d$ increases, the distinctions among trajectories within the same fiber become progressively less significant for the query.

An unstable trajectory space is one in which nearby histories \emph{diverge indefinitely}: no finite compression makes the fiber query-flat.

The Stability Principle: a trajectory space is stable for a query family when the projection of nearby trajectories onto query-answer space contracts with increasing interaction order.

\section{The Three Historical Universes}

\textbf{Outcome A (exponential stability).} History becomes rapidly irrelevant. Compression is efficient; reconstruction is tractable. Science is easy; institutions are reformable.

\textbf{Outcome B (polynomial stability).} History remains important for a long time. Compression is possible but costly; reconstruction requires large depth.

\textbf{Outcome C (no stability).} History never becomes negligible. Compression fails; reconstruction is impossible; explanation itself breaks down.

\section{Stability and Explanation}

Why do explanations work? The framework's answer: explanations work because most historical influence eventually becomes negligible. A good explanation is a compression sufficient for the query families relevant to understanding and prediction. In a world without stability, nothing could be summarized, modeled, or reconstructed. The possibility of explanation is a consequence of stability.

\section{Query-Relativity of Stability}

Stability is not intrinsic to the trajectory space. It is a relational property of the trajectory space, the query family, and the interaction decomposition. A physical system may be exponentially stable for thermodynamic queries and chaotic for microscopic queries. This query-relativity prevents treating ``stable systems'' as an intrinsic category.

\textbf{Stability determines when compression is possible. Projection determines what is preserved.}
